%% =====================================================================
%%  T-27-02  The Eagle's Ice
%%  -------------------------------------------------------------------
%%  One idea per file. Core is mandatory. Slots in canonical order.
%%  Max 700 words. No layout commands. No filler. New source material
%%  for this entry goes in sources/B1/T-27-02.tex -- never by
%%  rewriting this file (docs/RULES.md R0.1).
%%  =====================================================================
%% @kind: topic
%% @id: T-27-02
%% @title: The Eagle's Ice
%% @subtitle: The selective-ruthlessness licence does not exist
%% @chapter: c27
%% @order: 20
%% @status: stable
%% @sources: B1
%% @updated: 2026-09-19

\topic{T-27-02}{The Eagle's Ice}{The selective-ruthlessness licence does not exist}

\begin{core}
The licence the method most often grants itself is selectivity: ruthless with strangers, decent with friends. The source's finding is that this licence is not on offer. Conduct is per-trait, not per-target --- exploit some people and you will exploit anyone --- because behaviour barely changes from one relationship to the next. The selecting dissolves, and the pattern generalises to everyone, including the friends it was supposed to spare.
\end{core}
\begin{mechanism}
Every ruthless operator in the source's gallery carried the same private plan --- never cheat a friend, be good to the people who are good to me --- and all failed the same way. The unit of conduct is the trait, not the target: a man who mistreats his enemies undercuts his friends too, because the habit does not check the name before it fires. The chapter's image: an eagle spots a fish frozen in a block of ice drifting toward a waterfall, dives to snatch it, sticks, and goes over the falls with the ice. The snatch that was meant to be the once-only exception is the moment the pattern locks. Behaviour has directional momentum --- each selfish act makes the next feel natural, each restraint makes the next restraint easier --- so the selective operator is not standing anywhere; he is mid-slide, calling the slide a strategy.
\end{mechanism}
\begin{conditions}
Strongest for habitual conduct --- the slide runs on repetition, not intent, so it shows in anyone's history: the manager who grinds subordinates is grinding his family in instalments he explains as exceptions. Visible and reversible early; nearly invisible to the operator mid-slide, who experiences each act as discrete. Weakened only by institutional partition --- roles that genuinely fence conduct (the judge off the bench) --- and private partitions are just the slide wearing a costume.
\end{conditions}
\begin{application}
  \item Audit by trait, not incident: list how you treat everyone you have power over. The list should converge on one person.
  \item Treat the thought \emph{just this once} as the alarm. That thought is the dive, not the exception.
  \item To keep decency toward friends, practise decency where no friend is watching. The trait is trained in private.
  \item When you catch the slide, correct structurally --- change the incentives you live under --- not by resolve, which the slide outlasts.
  \item Extend the audit to the good direction: reliability also generalises. Build the trait you want repeated.
\end{application}
\begin{feedback}
  \item Working: the audit finds the same person everywhere. Boring, and safe.
  \item Working: strangers and intimates describe you with the same adjectives. The trait is one thing.
  \item Failing: your exceptions have exceptions. The pattern is locking.
  \item Failing: the people who know you best begin discounting your warmth the way strangers do. They have seen the trait switch on.
\end{feedback}
\begin{failure}
The law convicts in both directions. It convicts the selective operator --- but it also convicts the reformed one: because conduct changes little, repair is slow and the friends lost early stay lost; the audit's honesty does not shorten the sentence. And the same generality cuts for instrumental loyalty --- a person can be reliably decent for tool-maintenance reasons alone. The trait converges either way; only the ledger shows which.
\end{failure}
\begin{countermeasures}
  \item Evaluate operators by how they treat people who can do nothing for them. The trait performs where the audience is not.
  \item Listen for the ice: a counterpart describing their current mark with \emph{this one's different} is mid-dive.
  \item In your own house, the phrase \emph{they're family} should end a calculation, not open one. Calculations that start there are already over.
  \item Note who is warm to you and cold to waiters. You are watching the generalised trait rehearse --- and auditioning for a later scene.
\end{countermeasures}

\sources{B1}
\seesources
\seealso{T-27-01, T-06-03, T-28-01}
